\section{矩阵的四个基本子空间：输入、输出与丢失的信息}
\label{sec:03-Linear-Algebra/02-Vector-Spaces/03}

一份回归任务里，你不小心把身高放进去两次：一列以厘米记，一列以英寸记。模型照常拟合，预测值也挑不出毛病，可每换一个随机种子，这两列的系数就换一副面孔——一个 \(+3.2\) 配一个 \(-1.2\)，下次变成 \(-0.9\) 配 \(1.8\)，加权之后完全一样。你想解释``身高每增加一厘米，预测值上升多少''，却发现这个问句没有唯一答案。

问题不出在优化器，出在矩阵。两列之间存在一个恒等式，沿着它挪动参数，输出纹丝不动。要把这类现象说清楚，需要对一个矩阵同时问四个问题：输出端能到达哪里？哪些输出永远够不着？输入端的哪些方向被送了出去？哪些方向被就地抹掉？四个问题各对应一个子空间，它们合起来构成读懂任何矩阵的最小工具箱。

\subsection{输出端：能到达的地方与够不着的地方}

矩阵乘向量是把列按系数组合：\(Ax=x_1a_1+\cdots+x_na_n\)。所以当 \(x\) 跑遍 \(\R^n\)，\(Ax\) 跑遍的恰好是列张成的集合，即\emph{列空间} \(\operatorname{Col}(A)\subseteq\R^m\)。可解性由此变成一句几何判断：

\[
Ax=b\ \text{有解}\quad\Longleftrightarrow\quad b\in\operatorname{Col}(A).
\]

方程组有没有解，与方程个数和未知数个数的大小关系无关，只看右端向量落不落在列的触及范围里。消元中那一行 \([0\;\cdots\;0\mid c]\)、\(c\neq0\)，就是在报告 \(b\) 多出了一个列空间够不到的分量。列空间的维数称为矩阵的\emph{秩} \(\rank(A)\)，等于主元列的个数。

输出空间里还有一块列空间碰不到的区域。若某个 \(y\in\R^m\) 与 \(A\) 的每一列都正交，即 \(A^{\trans}y=0\)，那么对任意 \(x\) 都有 \(y^{\trans}(Ax)=0\)：无论输入是什么，输出在 \(y\) 方向上的分量恒为零。这些 \(y\) 构成\emph{左零空间} \(\mathcal N(A^{\trans})\)，名字里的``左''来自等价写法 \(y^{\trans}A=0\)。它是理想读数必须满足的约束集合：拿到一批实测数据，若 \(y^{\trans}b\) 明显不为零，说明数据与模型之间存在无法调和的矛盾——噪声、标定漂移，或者方程本身漏了项。

\subsection{输入端：被送出去的方向与被抹掉的方向}

输入空间的划分同样是两块。\emph{零空间}

\[
\mathcal N(A)=\{x\in\R^n:Ax=0\}
\]

收集所有被压成零的输入。它的分量为什么难缠，看一眼就明白：若 \(z\in\mathcal N(A)\)，则对任意 \(x\) 有 \(A(x+z)=Ax\)。从 \(A\) 的视角，\(x\) 与 \(x+z\) 是同一个东西。开头那两列身高，差的正是这么一个 \(z\)——两列之间的固定换算关系让某个系数组合落进了零空间，于是参数不可辨识，多解彼此等价。这不是数值精度问题，再好的求解器也分不清 \(A\) 本来就不区分的两个输入。

剩下的那一块是\emph{行空间} \(\operatorname{Row}(A)\subseteq\R^n\)，由 \(A\) 的行张成。行空间与零空间的关系有一行就能说完的理由：\(Ax=0\) 逐行展开，就是说 \(x\) 与每一行的内积都为零，于是与行的任何组合的内积也为零。换句话说，零空间里的向量与行空间里的向量两两正交。同样的推理搬到 \(A^{\trans}\) 上，得到列空间与左零空间正交。

\begin{center}
\begin{tikzpicture}[x=1cm,y=1cm,font=\small,>=stealth]
  \draw[black!45] (0,-2.7) rectangle (3.4,2.7);
  \node[above,font=\footnotesize] at (1.7,2.72) {输入空间 \(\R^n\)};
  \draw[black!55,fill=blue!12] (0.25,0.35) rectangle (3.15,2.4);
  \node[font=\footnotesize] at (1.7,1.72) {行空间};
  \node[font=\footnotesize] at (1.7,1.05) {\(\dim=r\)};
  \draw[black!55,fill=black!6] (0.25,-2.4) rectangle (3.15,-0.35);
  \node[font=\footnotesize] at (1.7,-1.05) {零空间};
  \node[font=\footnotesize] at (1.7,-1.72) {\(\dim=n-r\)};
  \node[font=\footnotesize] at (0.7,0) {\(\perp\)};
  \begin{scope}[shift={(7.2,0)}]
    \draw[black!45] (0,-2.7) rectangle (3.4,2.7);
    \node[above,font=\footnotesize] at (1.7,2.72) {输出空间 \(\R^m\)};
    \draw[black!55,fill=blue!12] (0.25,0.35) rectangle (3.15,2.4);
    \node[font=\footnotesize] at (1.7,1.72) {列空间};
    \node[font=\footnotesize] at (1.7,1.05) {\(\dim=r\)};
    \draw[black!55,fill=black!6] (0.25,-2.4) rectangle (3.15,-0.35);
    \node[font=\footnotesize] at (1.7,-1.05) {左零空间};
    \node[font=\footnotesize] at (1.7,-1.72) {\(\dim=m-r\)};
    \node[font=\footnotesize] at (2.7,0) {\(\perp\)};
    \fill (0.15,0) circle (1.8pt);
    \node[font=\footnotesize] at (0.62,0) {\(0\)};
  \end{scope}
  \draw[->,thick] (3.5,1.4) -- (7.1,1.4);
  \node[above,font=\footnotesize] at (5.3,1.45) {\(A\)：一一对应};
  \draw[->,thick,black!60] (3.5,-1.5) -- (7.25,-0.1);
  \node[below,font=\footnotesize] at (5.3,-0.95) {\(A\)：整块压到零};
\end{tikzpicture}

\emph{一个矩阵把两个空间各劈成互相垂直的两半。上半层的维数两边都是 \(r\)：\(A\) 在行空间与列空间之间建立一一对应。下半层是两端的盲区——输入里的零空间送不出任何信号，输出里的左零空间收不到任何信号。四个子空间和一个数 \(r\)，构成矩阵的完整账目。}
\end{center}

\begin{remark}
正交是这里最省力的说法，但严格谈论垂直需要内积，本单元的公理里还没有它。上面那句``两两正交''用的是 \(\R^n\) 中现成的点积，直和分解 \(\R^n=\operatorname{Row}(A)\oplus\mathcal N(A)\) 与 \(\R^m=\operatorname{Col}(A)\oplus\mathcal N(A^{\trans})\) 的完整论证留给 \hyperref[sec:03-Linear-Algebra/03-Orthogonality/01]{``正交与正交补''}。眼下只需接受这张图的形状。
\end{remark}

\subsection{秩\texorpdfstring{--}{-}零度定理：自由度守恒}

四个维数不是彼此独立的四个数。消元把 \(n\) 个变量分成 \(r\) 个主元变量和 \(n-r\) 个自由变量，每个自由变量恰好贡献一个独立的零空间方向，于是零空间维数是 \(n-r\)。

\begin{theorem}[秩--零度定理]
设 \(A\) 有 \(n\) 列，则 \(\rank(A)+\dim\mathcal N(A)=n\)。更一般地，对有限维空间 \(V\) 上的线性变换 \(T:V\to W\)，\(\dim V=\dim\ker T+\dim\operatorname{im}T\)。
\end{theorem}

证明骨架：在核里取一组基，把它扩充成整个 \(V\) 的基（上一节的扩充操作保证这一步可行）；扩充进来的那些向量在 \(T\) 作用下的像张成整个 \(\operatorname{im}T\)，且线性无关——若它们的某个组合被 \(T\) 送到零，这个组合就落在核里，于是能用核的基表出，与整体线性无关矛盾。数一数三组向量的个数，等式就出来了。

值得注意的是这条定理要什么、不要什么。它只要求\emph{输入空间有限维}；对输出空间的维数、对 \(A\) 的结构、对系数域，一概不作要求。反过来说，输入无限维时它确实会塌。取全体数列构成的空间，右移算子把 \((a_1,a_2,\ldots)\) 送成 \((0,a_1,a_2,\ldots)\)，核只有零元，像却是真子空间——``维数相减''这个动作在这里根本没有定义。

读法上，这个等式不是列数的重新分配，而是一份自由度账单：输入端共 \(n\) 个独立方向，\(r\) 个活着到达输出端，另外 \(n-r\) 个死在半路。想让一个变换少丢信息，只能提高 \(r\)，没有别的办法。这也解释了开头的现象——特征列里塞进一个线性可推导的重复列，\(n\) 增加了 \(1\)，\(r\) 却不变，零空间白白多出一维，多出来的自由度全部流向``改不动预测值的参数摆动''。

\subsection{最小范数解为什么落在行空间里}

把上面的分解用一次，能得到一条很实用的结论。任意输入唯一拆成 \(x=x_{\mathrm{row}}+x_{\mathrm{null}}\)，于是

\[
Ax=Ax_{\mathrm{row}}+Ax_{\mathrm{null}}=Ax_{\mathrm{row}}.
\]

零空间分量在输出里蒸发。所以当 \(Ax=b\) 相容时，全体解构成仿射集 \(x_p+\mathcal N(A)\)，是一条（或一片）平行于零空间的集合；沿着它移动不改变输出，只改变解本身的长度。想在这堆等价的解里挑一个，最自然的挑法是挑最短的那个——而最短的那个必定不含任何零空间分量，也就是落在行空间里。

\begin{center}
\begin{tikzpicture}[x=1cm,y=1cm,font=\small,>=stealth]
  \draw[black!35,->] (-0.8,0) -- (3.8,0);
  \draw[black!35,->] (0,-2.0) -- (0,3.5);
  \draw[blue!65!black,thick] (-0.3,-0.3) -- (3.1,3.1);
  \node[font=\footnotesize,blue!55!black] at (2.75,2.55) {行空间};
  \draw[black!55,dashed] (0,0) -- (1.7,-1.7);
  \node[font=\footnotesize,black!65] at (1.95,-1.55) {零空间};
  \draw[black!75,thick] (0.1,2.9) -- (3.3,-0.3);
  \node[font=\footnotesize] at (0.55,3.05) {解集 \(x_p+\mathcal N(A)\)};
  \draw[->,very thick,blue!65!black] (0,0) -- (1.5,1.5);
  \fill (1.5,1.5) circle (1.6pt);
  \node[font=\footnotesize,above left] at (1.5,1.5) {\(x_+\)};
  \draw[->,black!70] (0,0) -- (2.7,0.3);
  \fill (2.7,0.3) circle (1.6pt);
  \node[font=\footnotesize,below] at (2.75,0.25) {\(x_1\)};
\end{tikzpicture}

\emph{相容方程的解集是一条平行于零空间的直线，线上每个点给出同一个 \(Ax\)。它与行空间恰好交于一点 \(x_+\)，那是全体解中最短的一个；任何别的解 \(x_1\) 都在 \(x_+\) 上叠了一段零空间位移，只增长度、不改输出。}
\end{center}

这就是伪逆给出的那个解，也是``最小范数解''这一说法的几何来源。欠定问题里之所以默认取它，并不是因为它更接近真值，而是因为需要一个不依赖随机初值的确定挑法。挑法的代价与替代方案（例如改用正则化把偏好写进目标函数）在 \hyperref[sec:03-Linear-Algebra/07-SVD-and-Data-Applications/02]{``伪逆与病态问题''}里详细算账。

\subsection{一个矩阵的完整账目}

拿

\[
A=\begin{pmatrix}1&2&3\\2&4&6\end{pmatrix}
\]

走一遍。第二行是第一行的两倍，三列都指向 \((1,2)^{\trans}\)，所以 \(r=1\)。列空间是 \(\R^2\) 中由 \((1,2)^{\trans}\) 张成的直线；行空间是 \(\R^3\) 中由 \((1,2,3)^{\trans}\) 张成的直线；零空间是平面 \(x_1+2x_2+3x_3=0\)，维数 \(3-1=2\)；左零空间由 \((-2,1)^{\trans}\) 张成，维数 \(2-1=1\)。输入端 \(1+2=3\)，输出端 \(1+1=2\)，两侧都闭合。三列输入方向里只有一条通道能把信息送到输出，其余两维原地消失。

四组基全部来自同一次消元，取法只需记住行操作的不对称：RREF 的非零行直接给出行空间的基；主元位置对应的\emph{原矩阵}的列给出列空间的基；令自由变量轮流取 \(1\)、其余取 \(0\) 回代，给出零空间的基；对 \(A^{\trans}\) 重复一遍，给出左零空间的基。

维数账目甚至不需要看到矩阵内容。若 \(A\) 是 \(5\times 8\) 且 \(\rank(A)=3\)，立刻可知 \(\dim\operatorname{Col}(A)=\dim\operatorname{Row}(A)=3\)、\(\dim\mathcal N(A)=5\)、\(\dim\mathcal N(A^{\trans})=2\)。于是两个判断当场成立：\(Ax=b\) 不可能对所有 \(b\in\R^5\) 有解，因为列空间只占三维；相容时也不可能唯一，因为解集至少是五维的一大片。矩阵的尺寸只是门面，秩才是它的吞吐量。

\subsection{记住这张图，忘掉计算}

输入被切成行空间与零空间，输出被切成列空间与左零空间，中间只有 \(r\) 维能穿过去。存在性由列空间管，唯一性由零空间管，冗余由 \(n-r\) 计量，数据自洽性由左零空间检验——四个老问题，一张图一次回答。

同一张图会在后面反复出现，只是名字换掉。测量里它是可分辨性与一致性检验；模型辨识里它是参数能不能从数据中恢复；到了 \hyperref[sec:03-Linear-Algebra/07-SVD-and-Data-Applications/01]{``任意矩阵的奇异值分解''}，它会被加工到极精确的形态：在两端各选一组正交基之后，\(A\) 的动作简化为 \(r\) 条独立的缩放通道，缩放倍数为零的方向就是这里的盲区。

图里还留了一个没说破的地方。行空间到列空间的那条``一一对应''，是矩阵 \(A\) 在标准基下的行为；换一组基去看同一个变换，矩阵里的数字会全部改写，可这四个子空间的维数一个都不会动。下一节就把变换与它的矩阵分开，看清哪些东西属于操作本身，哪些只是记账方式。
